\section{Sailing North}

Start the adventure with the following:

\quoteBox{
	Your ship dashes through the waves as you are sailing north. The windspren are chasing your vessel and the anticipationspren are streaming around you - what awaits you behind the horizon?
}

\subsection{Plot Hook}

Unless you started the adventure with "\nameref{chap:tutorial}", add the following to the introduction:

\quoteBox{
	You embarked on this ship with a goal to find a lost Thaylen Expedition. All contact with them was lost several days ago and you are now racing time to save those poor souls from any dangers that may lurk on the islands of the Reshi Sea. And perhaps - to earn some spheres in the process.
}

\reasonBox{
	If you are running this adventure alongside some bigger campaign, you can introduce your own plot hooks instead. Perhaps the characters are already sailing on the Reshi Sea for some other reason and they will encounter the lost expedition by accident, starting a side quest spanning for 5 character levels? You can adjust the reasons for the party to travel to your taste and the backstory of the party.
}

\subsection{Search for the Island}

The characters must find tracks of the lost expedition. You can run this as an endeavor. The search may take several days - consider presenting this endeavor as a montage of small scenes happening during the travel. Possible approaches to the endeavor may include:

\scriptedOption{Reading the map.}{
	One of the characters may have a map with marked expedition's path. Using it requires \skillCheck{10}{Deduction} to properly determine where the expedition might have been lost.
}

\scriptedOption{Ship navigation.}{
	One must know their way with the ship to properly sail on the seas. \skillCheck{15}{Lore} can be used to determine if a character knows how to work with the ship.
}

\scriptedOption{Reading the wind and waves.}{
	The sea can be unpredictable. If a character knows how to read its mood, they can attempt \skillCheck{12}{Survival} to help navigate to safety.
}

\scriptedOption{Hiding from the highstorm.}{
	Though this far west the storms are not as devastating, it is still good idea to avoid them. Use \skillCheck{10}{Survival} to check if a character successfully avoided the highstorm.
}

\scriptedOption{Asking local people.}{
	Perhaps someone saw the expedition ship? Use \skillCheck{7}{Persuasion} to verify if you can get some information from the nearby Reshi Islands
}

\scriptedOption{Spotting the ship.}{
	A character can attempt \skillCheck{12}{Perception} to see if they can spot the expedition ship.
}

\scriptedOption{Pure luck.}{
	Sometimes all you can count on is your luck. A character may roll d20, a result 10 or higher means they were lucky enough with their search.
}

\subsection{Opportunities And Complications}

\oppCompTable{
	\oppCompRow{\iconOpp[15pt]}{You managed to catch some delicious fish for dinner.}
	\oppCompRow{\iconComp[15pt]}{The sea is not for you. You got sea sick.}
}

\creditedArtFigure{width=0.45\textwidth}
					{images/island.png}
					{AI}
					{Approaching \theIsland{} Reshi Isle.}
					{fig:island}


\subsection{\pursuitResolveHeader}

\pursuitReq{6}{3}

If the party fails during the endeavor they will be still able to find the lost expedition eventually. However, they will become \exhausted{} until they are able to rest on the island.

\reasonBox{
	This part of the adventure is rather short, despite taking more time for the characters. If you wish, you can expand this scene, allowing for more roleplay, character interaction, and/or adding some combat.   
}


